\documentclass[11pt]{article}
\usepackage[margin=1in]{geometry}
\usepackage{amsmath,amssymb}
\usepackage{booktabs}
\usepackage{array}
\usepackage{enumitem}
\usepackage{hyperref}

\title{Project Proposal: Bayesian Switching Linear Dynamical Systems\\
\large Regime-Aware Modeling for Financial Time Series}
\author{Course: AI61004 \quad \texttt{GROUP 10}}

\begin{document}
\maketitle

\section*{Group Members}
\begin{table}[h!]
\centering
\begin{tabular}{p{8cm} p{4cm}}
\toprule
\textbf{Name} & \textbf{Roll Number} \\
\midrule
Harshwardhan Walunjkar & 23MA10077 \\
Krish Shekhar & 23MA10029 \\
Nishad Gaikwad & 23MA10038 \\
Prithviraj Solanke & 23IE10048 \\
Siddharth Nema & 23MA10078 \\
\bottomrule
\end{tabular}
\end{table}

\section{Objective}
The objective of this project is to develop a Bayesian regime-switching framework for financial time series that can:
\begin{itemize}[leftmargin=1.5em]
    \item infer latent market regimes (e.g., calm, stressed, crisis),
    \item model within-regime dynamics using linear state-space components,
    \item capture regime persistence and state-dependent switching behavior,
    \item provide actionable outputs for risk monitoring and strategy design.
\end{itemize}
The project is motivated by the need for models that are both probabilistically principled and practically useful in quantitative finance.

\section{Problem Formulation}
Given observations $y_{1:T}$ (returns and market covariates), we assume a latent discrete state sequence $z_{1:T}$ and a latent continuous state sequence $x_{1:T}$:
\begin{align}
z_t \mid z_{t-1} &\sim \pi_{z_{t-1}},\\
x_t &= A^{(z_t)}x_{t-1} + e_t^{(z_t)}, \quad e_t^{(z_t)} \sim \mathcal{N}(0,Q^{(z_t)}),\\
y_t &= Cx_t + w_t, \quad w_t \sim \mathcal{N}(0,R).
\end{align}

To avoid fixing the number of regimes, we use a sticky HDP prior:
\begin{align}
\beta &\sim \mathrm{Dir}\!\left(\tfrac{\gamma}{L},\dots,\tfrac{\gamma}{L}\right),\\
\pi_j &\sim \mathrm{Dir}(\alpha\beta + \kappa \delta_j),
\end{align}
where $\kappa$ encourages persistent states.

To model non-geometric, state-dependent switches, we additionally use recurrent transitions:
\begin{align}
z_{t+1}\mid z_t,x_t \sim \pi_{\mathrm{SB}}(R_{z_t}x_t + r_{z_t}),
\end{align}
with stick-breaking logistic probabilities and P{\'o}lya--Gamma augmentation for Bayesian inference.

\section{Planned Models}
\begin{enumerate}[leftmargin=1.5em]
    \item \textbf{Sticky HDP-AR-HMM} (switching autoregressive model with nonparametric regime count).
    \item \textbf{Sticky HDP-SLDS} (full switching linear dynamical system with latent continuous states).
    \item \textbf{Recurrent switching model} with state-dependent transitions (rSLDS/rAR-HMM formulation).
\end{enumerate}
The three models will be compared under a common real-data evaluation protocol.

\section{Data Requirements and Preparation}
\begin{itemize}[leftmargin=1.5em]
    \item Primary series: broad equity index returns (e.g., S\&P 500 / NIFTY 50).
    \item Optional covariates: VIX changes, yield-curve spread, credit spread, volume/liquidity proxies.
    \item Frequency: daily; history target: at least 10--20 years where available.
\end{itemize}



\section*{References}
\begin{enumerate}[leftmargin=1.5em]
    \item Fox, E. B., Sudderth, E. B., Jordan, M. I., \& Willsky, A. S. (2008). \emph{Nonparametric Bayesian Learning of Switching Linear Dynamical Systems}. NeurIPS.
    \item Linderman, S. W., Johnson, M. J., Miller, A. C., Adams, R. P., Blei, D. M., \& Paninski, L. (2017). \emph{Bayesian Learning and Inference in Recurrent Switching Linear Dynamical Systems}. AISTATS.
    \item Linderman, S. W. et al. (2017). \emph{Supplementary Material: Bayesian Learning and Inference in Recurrent Switching Linear Dynamical Systems}.
    \item Teh, Y. W., Jordan, M. I., Beal, M. J., \& Blei, D. M. (2006). \emph{Hierarchical Dirichlet Processes}. JASA.
\end{enumerate}

\end{document}

